\section{Related Work}
\label{sec:related}

\paragraph{Representations of units larger than a token.}
The closest earlier work is \citet{kaplan-etal-2025-tokens}, who show
that when a word is split into several subword tokens, early attention
aggregates the pieces into the final subtoken and feed-forward layers then write
a whole-word representation into that position. They read these states out with
Patchscopes and find that the model can regenerate words for which no single
vocabulary token exists, which they interpret as a latent lexicon broader than
the tokenizer's vocabulary. Their unit is a word the model has finished reading; ours is a phrase it has
not. The two share a method but not a claim: pulling together evidence already
seen is a weaker thing than anticipating evidence not yet seen.
\citet{reif-etal-2026-vocab} build on the same picture to reshape model
vocabularies, and \citet{artzy-schwartz-2024-attend} add evidence that attention
gathers information early and later layers work on the result.

\paragraph{Idioms and memorised sequences.}
\citet{haviv-etal-2023-understanding} study how transformers recall memorised
idiom completions, and find the process to be two-stage: early layers promote
the memorised token to the top of the output distribution, and upper layers
raise the model's confidence in it. They study how a single target token is recalled; we measure how far ahead
recall reaches and how it varies across phrases and sentences. Our idiom set is drawn from
MAGPIE \citep{haagsma-etal-2020-magpie} and
LIdioms \citep{moussallem-etal-2018-lidioms}.

\paragraph{Anticipating future tokens from one hidden state.}
Our question was posed, in a more general form, by Future
Lens \citep{pal-etal-2023-future}: given the hidden state of a single token at
position $t$, can the tokens at positions $\geq t+2$ be recovered? They answer
it with linear approximation and with causal intervention (transplanting the
state into a different context and reading what the model then produces), and
report better than $48\%$ accuracy at some layers of GPT-J-6B. That transplant
is the technique Patchscopes later generalised, and our Patchscopes readout is a
direct descendant of it.

We differ in what is anticipated and in what it is checked against. Future Lens
scores the recovered tokens against whatever the model would have written next;
we score them against the phrase, which is fixed in advance and does not move
with the model. That lets us repeat a phrase across many sentences
(\S\ref{sec:stability}) and set the transplant against the model's own writing
on matched instances (\S\ref{sec:disagreement}). Future Lens compares two
\emph{decoding} methods against the model's output; we compare a decoding
method against the model's behaviour, and find that they disagree in ways that
change conclusions.

\paragraph{Reading hidden states.}
The logit lens \citep{nostalgebraist-2020-logitlens} projects an intermediate
state through the unembedding to ask what token it would predict.
Patchscopes \citep{ghandeharioun-etal-2024-patchscopes} generalises both it and
the Future Lens intervention by injecting the state into an arbitrary prompt and
letting the model write it out, which lifts the restriction that the answer be
a single vocabulary item and so allows multi-token readouts. All three act on a representation at inference time, and all raise the same
question: is what comes out what the model would otherwise have used? We test
that directly in \S\ref{sec:disagreement}.

\paragraph{Probing internal states for future correctness.}
\citet{orgad-etal-2025-llms} show that hidden states encode more about the
truthfulness of a forthcoming answer than the output distribution reveals, that
the signal concentrates in particular tokens, and that probes trained on it
generalise poorly across datasets. Our probe is the same kind of thing pointed at a different question, not whether
an answer will be true but whether a phrase will be finished, and it behaves
similarly, including the limited transfer across categories in
Appendix~\ref{app:details}. In a similar spirit,
\citet{gekhman-etal-2024-fine} treat a model's pre-existing knowledge of an item
as a measurable quantity predicting downstream behaviour.

\paragraph{Exploiting anticipation for efficiency.}
If a model's state already determines several upcoming tokens, that redundancy
can be spent. Speculative decoding \citep{leviathan-etal-2023-fast} drafts
several tokens with a cheaper process and verifies them with the target model;
multi-token prediction heads \citep{gloeckle-etal-2024-better} train the model
to emit several tokens at once; and LayerSkip
\citep{elhoushi-etal-2024-layerskip} exits early and self-speculates.
\citet{aynetdinov-akbik-2025-babies} show that a multi-token prediction
objective helps even 130M-parameter models under severe data constraints,
suggesting the capacity to train for lookahead is not confined to large models.

Closest to the use our probe suggests is CLP \citep{xie-zhou-2026-clp}, which
trains one small linear layer on the hidden state to set how many drafted tokens
to accept at each step, reporting $1.14\times$--$1.29\times$ on Qwen2.5 without
measurable quality loss: the same kind of read as our probe, used to make a
decoding decision rather than to answer a question. All of these methods assume
useful anticipation exists. We measure it directly, for one well-defined class
of continuations, and say when it is reliable. We implement no decoding
method here and report no speed-ups of our own.
